\begin{lemma}[Finite targets kill the torsion from a group above a band]
\label{B:group-band-torsion-killed}
Let $H$ be a finite group acting on finite nonempty row and column
sets $I,J$, and let $S=H\amalg(I\times J)$ be the corresponding
group-band monoid. Every map $BS\to X$ to a connected finite CW
complex kills the torsion subgroup of $H_2(BS;\Z)$.
\end{lemma}
\begin{proof}
Theorem~\ref{E:group-band-cofiber} identifies $BS$ with the homotopy
cofiber of $(I*J)_{hH}\to BH$. The Borel construction is a finite
graph of classifying spaces of finite stabilizer groups: its vertex
groups are stabilizers of row and column orbits, and its edge groups
are stabilizers of pairs. Miller's full mapping-space theorem
\cite[Theorem~A]{miller1984} makes evaluation
$\operatorname{Map}(BL,X)\to X$ a weak equivalence for each of
these finite groups $L$. Naturality and the finite homotopy-colimit
description show that collapse to the orbit graph induces a weak
equivalence on mapping spaces into $X$. The same holds for
$BH\to *$. Applying the cofiber mapping-space sequence shows that
\[
 BS\longrightarrow\Sigma\bigl((I*J)/H\bigr)
\]
induces a weak equivalence on mapping spaces into $X$. The orbit
graph is finite and connected, so its suspension is a finite wedge
of two-spheres. Every map from $BS$ to $X$ therefore factors up to
homotopy through a space with torsion-free second homology, which
proves the assertion.
\end{proof}

\begin{theorem}[A three-by-three minimal band has no three-primary torsion]
\label{B:three-by-three-no-three}
Let $M$ be a finite monoid whose minimal ideal is a $3\times3$
rectangular band. If $BM$ has finite CW homotopy type, then
$H_2(BM;\Z)$ has no nonzero three-primary torsion.
\end{theorem}
\begin{proof}
Let $V,W$ be the row and column augmentation modules, each free
of rank two. Lemma~\ref{lem:B-band-Tor} identifies
$A=H_2(BM;\Z)$ with the balanced quotient of $W\otimes V$.
Suppose that $A$ has three-primary torsion. Choose
$\lambda:A\to\F_3$ nonzero on that torsion. In augmentation
bases, write its nonzero bilinear matrix as $L\in M_2(\F_3)$.
For the row and column transformations $\sigma_m,\tau_m$, let
$A_m,B_m$ be their integral augmentation matrices. Balance says
\begin{equation}\label{B:three-balance}
 B_m^{\mathsf T}L=LA_m\pmod3.
\end{equation}

We first record the complete small linear classification. Up to
relabeling the two three-point sets, transposition, and multiplying
$L$ by $-1$, its possibilities are the following five types:
\[
\begin{array}{c|c|c}
\text{type}&L&\text{number of projective matrices}\\ \hline
\text{two levels, two levels}&
 \left(\begin{smallmatrix}0&0\\0&1\end{smallmatrix}\right)&9\\
\text{two levels, three levels}&
 \left(\begin{smallmatrix}0&0\\1&2\end{smallmatrix}\right)&6\\
\text{three levels, three levels}&
 \left(\begin{smallmatrix}1&2\\2&1\end{smallmatrix}\right)&1\\
\text{rank two, first orbit}&
 \left(\begin{smallmatrix}0&1\\1&0\end{smallmatrix}\right)&18\\
\text{rank two, second orbit}&
 \left(\begin{smallmatrix}0&1\\1&1\end{smallmatrix}\right)&6.
\end{array}
\]
Here is a proof of completeness. The projective dual of the
three-point augmentation module has four points. One is represented
by values $(0,1,2)$ and is fixed by all permutations; the other three
are the two-level functions and are permuted transitively. Rank-one
matrices therefore give the first three orbits, of sizes $9,6,1$.
For rank two, the projective permutation action is the order-six
stabilizer of the invariant point in
$\operatorname{PGL}_2(\F_3)$, which has order $24$ and acts on four
points. Its left and right actions on projective isomorphisms have
two double cosets, according as the distinguished point maps to the
distinguished point or to one of the other three. Their sizes are
$6$ and $18$, represented by the two displayed invertible matrices.
The orbit sizes sum to $(3^4-1)/2=40$.

In the first, fourth and fifth types choose the displayed integral
matrix, whose entries are zero or one. Each side of a single entry
of \eqref{B:three-balance} is a difference of two entries of the
extended zero-one grid, with row and column zero adjoined. Their
difference lies between $-2$ and $2$. Divisibility by three therefore
forces integral equality. Hence $L$ defines an integral balanced
functional $A\to\Z$ lifting $\lambda$. Integral functionals kill
torsion, contrary to the choice of $\lambda$. Relabelings lift to
integral permutation maps, so the conclusion holds throughout each
of these three orbits.

Consider the mixed type. Write $L=b\otimes a$, with row values
$a=(0,1,2)$ and column values $b=(0,0,1)$. The rank-one balance
equation means that every element has a common slope
$\chi_m\in\{0,1,-1\}\subset\F_3$:
\[
 a_{\sigma_m(i)}=\chi_ma_i+c_m,\qquad
 b_{\tau_m(j)}=\chi_mb_j+d_m.
\]
Suppose an element with $\chi_m=1$ acts nontrivially on rows.
Its row action is a three-cycle. Its column action preserves both
$b$-fibers, because a nonzero translation cannot preserve a two-point
subset of $\F_3$. In the eventual cyclic local group of powers of
$m$, take its three-Sylow subgroup $H$, with identity $e$.
The row action remains transitive on all three rows. The column
permutation image has order at most two, so $H$ fixes every column
in $J'=\operatorname{im}\tau_e$. This image retains both $b$-values.
The corner contains
$S=H\amalg(\{0,1,2\}\times J')$. Its balanced presentation is
\[
 H_2(BS;\Z)
 \cong\ker(\Z J'\to\Z)\otimes
       \bigl(\ker(\Z^3\to\Z)\bigr)_{C_3}
 \cong(\Z/3)^{|J'|-1}.
\]
The functional $\lambda$ is nonzero on this group: take the
difference of columns with distinct $b$-values and the difference
of rows with distinct $a$-values. This contradicts
Lemma~\ref{B:group-band-torsion-killed} for the map to a finite CW
model of $BM$. For the nonunital corner inclusion, adjoin a fresh
identity to $S$; its old-identity corner has the same classifying
space and balanced presentation, exactly as in
Proposition~\ref{B:two-row-torsion}.

Consequently every slope-one row action is the identity. Any two
slope-minus-one row reflections have the same fixed point, since
two distinct reflections compose to a nontrivial slope-one
translation. Assign integral row values $-1,0,1$, placing zero at
that common fixed point and choosing the signs to recover $a$
modulo three up to an additive constant. If no reflection occurs,
choose these three lifts in any order with the same residues.
Keep the integral column values $(0,0,1)$. A slope-zero row action
is constant and a slope-zero column action has constant $b$-value;
slope-one actions preserve both functions; and a slope-minus-one
action negates the row function and sends $b$ to $1-b$.
Thus their integral linear parts have the same slope, and their
tensor product is an integral balanced lift of $\lambda$.
This again contradicts torsion detection.

It remains to treat the three-level, three-level type. Both value
functions are bijections to $\F_3$. A slope-zero pair consists of
two constant transformations. Every other pair consists of two
affine permutations with the same slope. The permutation pairs in
the image of $M$ form a finite group $P$. The preimage of $P$ in
$M$ is a finite submonoid mapping onto $P$, so a local maximal
group $H$ maps onto $P$ by Lemma~\ref{A:minimal-group}. Its identity
acts identically on the whole minimal band $K$. Hence the corner
contains $S=H\amalg K$.

The map $H_2(BS;\Z)\to H_2(BM;\Z)$ is an isomorphism. Indeed
the balanced presentations have exactly the same relations from
$P$, and every remaining action pair is constant on both sides,
so both augmentation actions are zero and it contributes no
relation. Naturality of Lemma~\ref{lem:B-band-Tor}, with the same
fresh-identity treatment of the corner, identifies this isomorphism
with the actual map of classifying spaces. But
Lemma~\ref{B:group-band-torsion-killed} says that this map kills
torsion. Thus $A$ is torsion-free in this last case, the final
contradiction.
\end{proof}
